\documentclass[11pt,aspectratio=169]{beamer}
\usetheme{SimplePlus}

\usepackage[T1]{fontenc}
\usepackage[utf8]{inputenc}
\usepackage{amsmath,amssymb,mathtools}
\usepackage{booktabs}
\usepackage{array}
\usepackage{appendixnumberbeamer}
\usepackage{hyperref}

\newcommand{\E}{\mathbb{E}}
\newcommand{\Prob}{\mathbb{P}}
\newcommand{\Var}{\mathrm{Var}}
\newcommand{\Cov}{\mathrm{Cov}}
\newcommand{\ATT}{\mathrm{ATT}}
\newcommand{\ATE}{\mathrm{ATE}}
\newcommand{\TOT}{\mathrm{TOT}}
\newcommand{\MMSE}{\mathrm{MMSE}}
\newcommand{\indep}{\perp\!\!\!\perp}
\newcommand{\argmin}{\mathop{\mathrm{arg\,min}}}
\newcommand{\derivbutton}[1]{\vspace{0.4em}\hfill\hyperlink{app:#1}{\beamerbutton{Appendix details}}}
\newcommand{\backbutton}[1]{\vspace{0.4em}\hfill\hyperlink{main:#1}{\beamerbutton{Back to main slide}}}

\title{Applied Econometrics: Session 2}
\subtitle{Making Regression Make Sense: Selection on Observables}
\author{PhD Intensive Course}
\institute{Lecture 2}
\date{\today}

\begin{document}

\begin{frame}
  \titlepage
\end{frame}



\begin{frame}{Why This Lecture Matters}
  Lecture 1 gave the causal language: potential outcomes, selection bias, and random assignment.
  \vspace{0.4em}

  Lecture 2 asks a harder question:
  \begin{block}{Core question}
    When we run a regression in observational data, what does that coefficient mechanically estimate, and when can we read it as causal?
  \end{block}

  \begin{itemize}
    \item Regression is always a statistical approximation device.
    \item Causality is an extra layer, not a free gift from software.
  \end{itemize}
\end{frame}

\section{Regression Fundamentals}

\begin{frame}{Regression as a Summary of the CEF}
  \begin{itemize}
    \item Start from the object we want to summarize: the conditional expectation function.
    \item Then define regression as a disciplined linear approximation to that object.
    \item Only after that do we ask whether the CEF itself is causal.
  \end{itemize}
\end{frame}

\begin{frame}{Why Regression Shows Up Everywhere}
  \begin{itemize}
    \item In experiments, regression with a treatment dummy reproduces treatment-control mean differences.
    \item In observational work, regression is a control strategy for observed confounders.
    \item In descriptive work, regression is a compact summary of how averages vary with covariates.
  \end{itemize}
  \begin{block}{Core point}
    Before arguing about causal interpretation, we should understand what regression always does mechanically.
  \end{block}
\end{frame}

\begin{frame}{The Conditional Expectation Function}
  \hypertarget{main:cef}{}
  For a scalar outcome \(Y_i\) and covariates \(X_i\), define
  \[
    m(x) \equiv \E[Y_i \mid X_i=x].
  \]
  \begin{itemize}
    \item The CEF is the average value of \(Y_i\) for units with the same covariates.
    \item If \(X_i\) is binary, the CEF is just two cell means.
    \item If \(X_i\) is continuous, the CEF is a function, often nonlinear.
    \item The CEF is about averages, not automatically about causality.
  \end{itemize}
  \begin{block}{Example}
    \(m(s)=\E[\log w_i \mid S_i=s]\) is the average log wage of workers with \(s\) years of schooling.
  \end{block}
\end{frame}

\begin{frame}{Law of Iterated Expectations}
  \hypertarget{main:lie}{}
  \[
    \E[Y_i]=\E\{\E[Y_i \mid X_i]\}.
  \]
  \begin{itemize}
    \item First average within covariate groups.
    \item Then average those group means using the distribution of \(X_i\).
    \item This identity is the bridge between cell-specific effects and population averages.
  \end{itemize}
  \begin{block}{Intuition}
    An unconditional average is just a weighted average of conditional averages.
  \end{block}
  \derivbutton{cef-foundations}
\end{frame}

\begin{frame}{CEF Decomposition Property}
  \hypertarget{main:cef-decomp}{}
  \[
    Y_i = \E[Y_i \mid X_i] + \varepsilon_i,
    \qquad
    \E[\varepsilon_i \mid X_i]=0.
  \]
  \begin{itemize}
    \item The first term is the part explained by \(X_i\).
    \item The residual is mean-independent of \(X_i\).
    \item Therefore the residual is uncorrelated with any function of \(X_i\).
  \end{itemize}
  \begin{block}{Why this matters}
    The CEF is not just descriptive; it is the unique decomposition of outcomes into an explained part and an orthogonal leftover.
  \end{block}
  \derivbutton{cef-foundations}
\end{frame}

\begin{frame}{The CEF Is the Best Unrestricted Predictor}
  \[
    \E[Y_i \mid X_i]
    =
    \argmin_{m(X_i)} \E\left[\left(Y_i-m(X_i)\right)^2\right].
  \]
  \begin{itemize}
    \item If you may use any function of \(X_i\), the CEF minimizes mean squared prediction error.
    \item This is an \(\MMSE\) statement, not a causal statement.
    \item The same logic explains why averages are central objects in econometrics.
  \end{itemize}
  \begin{block}{Takeaway}
    The CEF is the best possible predictor given the information in \(X_i\), but it may still be nonlinear and causally meaningless.
  \end{block}
\end{frame}

\begin{frame}{Population Regression Is a Linear Approximation Problem}
  \hypertarget{main:pop-reg}{}
  Define the population regression coefficient vector by
  \[
    \beta
    =
    \argmin_b \E\left[\left(Y_i-X_i'b\right)^2\right].
  \]
  First-order conditions imply
  \[
    \beta = \E[X_iX_i']^{-1}\E[X_iY_i].
  \]
  If \(e_i \equiv Y_i-X_i'\beta\), then
  \[
    \E[X_i e_i]=0.
  \]
  \begin{itemize}
    \item The regression residual is orthogonal to the regressors by construction.
    \item This orthogonality does not by itself imply causality.
  \end{itemize}
\end{frame}

\begin{frame}{Regression Anatomy}
  \hypertarget{main:anatomy}{}
  In the simple regression of \(Y_i\) on a constant and \(x_i\),
  \[
    \beta_1=\frac{\Cov(Y_i,x_i)}{\Var(x_i)}.
  \]
  In a multivariate regression, the coefficient on \(x_{ki}\) is
  \[
    \beta_k=\frac{\Cov(Y_i,\tilde{x}_{ki})}{\Var(\tilde{x}_{ki})},
  \]
  where \(\tilde{x}_{ki}\) is the residual from regressing \(x_{ki}\) on all other covariates.
  \begin{itemize}
    \item Every multiple-regression coefficient is a bivariate slope after partialling out other regressors.
    \item This is the algebraic core of Frisch-Waugh-Lovell intuition.
  \end{itemize}
  \derivbutton{anatomy-proof}
\end{frame}

\begin{frame}{Three Reasons Regression Is Useful}
  \hypertarget{main:reg-just}{}
  \begin{enumerate}
    \item If the CEF is linear, regression recovers it exactly.
    \item Even if the CEF is nonlinear, regression is the best linear predictor of \(Y_i\).
    \item Most importantly, regression is the best linear approximation to the CEF itself.
  \end{enumerate}
  \[
    \beta
    =
    \argmin_b
    \E\left[\left(\E[Y_i \mid X_i]-X_i'b\right)^2\right].
  \]
  \begin{block}{Preferred interpretation}
    Regression is best thought of as approximating the CEF, not as solving a philosophical problem about individual prediction.
  \end{block}
  \derivbutton{reg-justification}
\end{frame}

\begin{frame}{Saturated Models and Interactions}
  \begin{itemize}
    \item If all covariates are discrete and we include a dummy for every covariate cell, regression is saturated.
    \item Saturated regression reproduces the cell means exactly.
    \item Main effects without interactions impose additivity across covariates.
    \item Interactions matter whenever the effect of one covariate depends on another.
  \end{itemize}
  \begin{block}{Student warning}
    ``Controlling for \(X\)'' is not one thing. Controlling flexibly with interactions is very different from forcing a linear or additive specification.
  \end{block}
\end{frame}

\begin{frame}{Population First, Sample Second}
  \begin{itemize}
    \item The population coefficient \(\beta\) is the object of interest.
    \item The sample estimator is \(\hat{\beta}=(X'X)^{-1}X'Y\).
    \item Under random sampling and standard regularity conditions, \(\hat{\beta}\) converges to \(\beta\).
    \item Heteroskedasticity affects precision formulas, not the meaning of \(\beta\).
  \end{itemize}
  \begin{block}{Course sequencing}
    We will discuss inference in more detail later; today the priority is what the estimand means.
  \end{block}
\end{frame}

\section{Regression and Causality}

\begin{frame}{When Does Regression Become Causal?}
  \begin{itemize}
    \item Regression always approximates a CEF.
    \item But only a causal CEF can justify causal language for the coefficient.
    \item The key assumption is conditional independence, also called selection on observables.
  \end{itemize}
\end{frame}

\begin{frame}{From Predictive CEFs to Causal CEFs}
  \begin{itemize}
    \item A predictive CEF summarizes average outcomes across observed groups.
    \item A causal CEF summarizes average potential outcomes for a fixed reference population.
    \item To move from one to the other, we need a design or assumption that removes selection bias.
  \end{itemize}
  \begin{block}{Translation rule}
    Regression inherits its interpretation from the CEF it approximates. If the CEF is causal, regression is causal. If the CEF is not causal, regression is not causal.
  \end{block}
\end{frame}

\begin{frame}{Selection Bias Returns in Observational Data}
  \hypertarget{main:cia}{}
  Let \(C_i\) denote college attendance. Then
  \[
  \E[Y_i \mid C_i=1]-\E[Y_i \mid C_i=0]
  =
  \underbrace{\E[Y_{1i}-Y_{0i} \mid C_i=1]}_{\ATT}
  +
  \underbrace{\E[Y_{0i} \mid C_i=1]-\E[Y_{0i} \mid C_i=0]}_{\text{selection bias}}.
  \]
  \begin{itemize}
    \item The observed difference combines a causal term and a non-causal baseline term.
    \item In schooling applications, the bias is often thought to be positive because college-goers would earn more even without college.
  \end{itemize}
  \derivbutton{cia-proof}
\end{frame}

\begin{frame}{The Conditional Independence Assumption}
  \[
    \{Y_{0i},Y_{1i}\}\indep C_i \mid X_i.
  \]
  \begin{itemize}
    \item Conditional on observed covariates \(X_i\), treatment is as good as random.
    \item Equivalently, once we hold \(X_i\) fixed, treated and untreated units are comparable in potential outcomes.
    \item This is why the assumption is called selection on observables.
  \end{itemize}
  \begin{block}{Important}
    CIA is not testable from the observed outcome alone. It is a research-design claim about the treatment assignment process.
  \end{block}
\end{frame}

\begin{frame}{What CIA Buys Us}
  Under CIA,
  \[
    \E[Y_i \mid X_i,C_i=1]-\E[Y_i \mid X_i,C_i=0]
    =
    \E[Y_{1i}-Y_{0i} \mid X_i].
  \]
  \begin{itemize}
    \item Within an \(X_i\)-cell, treated and untreated units can be compared causally.
    \item Selection bias vanishes after conditioning on \(X_i\).
    \item The price is strong: every confounder that matters must be observed and correctly controlled.
  \end{itemize}
\end{frame}

\begin{frame}{Multi-Valued Treatment and Potential Outcome Schedules}
  When treatment intensity is \(S_i\) rather than a binary indicator, write
  \[
    Y_{si}\equiv f_i(s).
  \]
  CIA becomes
  \[
    Y_{si}\indep S_i \mid X_i \quad \text{for all } s.
  \]
  Then
  \[
    \E[Y_i \mid X_i,S_i=s]-\E[Y_i \mid X_i,S_i=s-1]
    =
    \E[f_i(s)-f_i(s-1)\mid X_i].
  \]
  \begin{itemize}
    \item This is the logic behind causal interpretations of schooling regressions.
    \item The treatment effect may vary with both \(s\) and \(X_i\).
  \end{itemize}
\end{frame}

\begin{frame}{From \(X\)-Specific Effects to Population Effects}
  The law of iterated expectations turns cell-specific causal effects into summary effects:
  \[
    \ATE
    =
    \E\left\{\E[Y_i \mid X_i,S_i=s]-\E[Y_i \mid X_i,S_i=s-1]\right\}.
  \]
  The effect on those actually treated at \(s\) uses different weights:
  \[
    \E\left\{\E[Y_i \mid X_i,S_i=s]-\E[Y_i \mid X_i,S_i=s-1]\mid S_i=s\right\}.
  \]
  \begin{itemize}
    \item Same cell-specific effects.
    \item Different target populations.
    \item Different weights over \(X_i\).
  \end{itemize}
\end{frame}

\begin{frame}{A Linear Causal Model That Leads to Regression}
  Suppose potential outcomes satisfy
  \[
    f_i(s)=\alpha+\rho s+\eta_i.
  \]
  If
  \[
    \eta_i = X_i'\gamma + v_i
    \qquad \text{and CIA holds given } X_i,
  \]
  then the observed outcome obeys
  \[
    Y_i=\alpha+\rho S_i+X_i'\gamma+v_i,
  \]
  with \(v_i\) orthogonal to \(S_i\) and \(X_i\).
  \begin{block}{Interpretation}
    In this special case, the regression coefficient \(\rho\) is the causal return to one additional unit of treatment.
  \end{block}
  \derivbutton{cia-proof}
\end{frame}

\begin{frame}{What Counts as a Good Control?}
  Good controls are variables that:
  \begin{itemize}
    \item are determined before treatment,
    \item help explain both treatment assignment and outcomes,
    \item make the CIA more plausible,
    \item do not themselves lie on the causal path from treatment to outcome.
  \end{itemize}
  \begin{block}{Rule of thumb}
    Good controls are confounders or proxies for confounders. They are not consequences of treatment.
  \end{block}
\end{frame}

\begin{frame}{The Omitted Variables Bias Formula}
  \hypertarget{main:ovb}{}
  Suppose the long regression is
  \[
    Y_i=\alpha+\rho S_i+A_i'\gamma+\varepsilon_i,
  \]
  where \(A_i\) denotes omitted ability or family-background variables.
  Then the short-regression slope satisfies
  \[
    \frac{\Cov(Y_i,S_i)}{\Var(S_i)}
    =
    \rho+\gamma'\delta_{AS},
  \]
  where \(\delta_{AS}\) is the vector of coefficients from regressions of \(A_i\) on \(S_i\).
  \begin{block}{Memorize this}
    Short \(=\) long \(+\) effect of omitted \(\times\) correlation of omitted with included.
  \end{block}
  \derivbutton{ovb-proof}
\end{frame}

\begin{frame}{Sign Intuition for OVB}
  \begin{itemize}
    \item If omitted ability raises wages, then elements of \(\gamma\) are positive.
    \item If ability is positively correlated with schooling, then \(\delta_{AS}\) is positive.
    \item The short schooling regression overstates the causal return.
  \end{itemize}
  \begin{block}{But do not overlearn the sign}
    OVB is not always upward. If omitted factors are negatively correlated with schooling, bias can go the other way.
  \end{block}
\end{frame}

\begin{frame}{Illustrative Schooling Example}
  \begin{center}
    \begin{tabular}{@{}lc@{}}
      \toprule
      Specification & Schooling Coefficient \\
      \midrule
      No controls & 0.132 \\
      + Age dummies & 0.131 \\
      + Family background and demographics & 0.114 \\
      + AFQT score & 0.087 \\
      + Occupation dummies & 0.066 \\
      \bottomrule
    \end{tabular}
  \end{center}
  \begin{itemize}
    \item The drop from \(0.132\) to \(0.087\) is consistent with positive omitted-ability bias.
    \item The final drop after adding occupation dummies is harder to interpret.
  \end{itemize}
  \begin{block}{Why harder?}
    Occupation is likely affected by schooling, so it is not a clean pre-treatment control.
  \end{block}
\end{frame}

\begin{frame}{Bad Controls}
  \hypertarget{main:bad-control}{}
  A bad control is a variable that is itself an outcome of the treatment.
  \begin{itemize}
    \item Example: controlling for occupation in a regression of wages on college.
    \item College affects occupation.
    \item Conditioning on occupation therefore changes the composition of the comparison groups.
  \end{itemize}
  \[
    \E[Y_i \mid W_i=1,C_i=1]-\E[Y_i \mid W_i=1,C_i=0]
  \]
  is generally not a causal effect.
  \begin{block}{Verbal summary}
    We stop comparing like with like because the treatment changes who enters the conditioned-on group.
  \end{block}
  \derivbutton{bad-control-proof}
\end{frame}

\begin{frame}{Proxy Controls Can Also Go Bad}
  \begin{itemize}
    \item Suppose you cannot observe latent ability directly.
    \item You may be tempted to control for a post-treatment proxy, such as later occupation or post-schooling test scores.
    \item That proxy may partly absorb the very channel through which treatment works.
  \end{itemize}
  \begin{block}{Safe question}
    Was this variable already fixed before treatment was chosen?
  \end{block}
  \begin{block}{Unsafe question}
    Is this variable merely correlated with the outcome?
  \end{block}
\end{frame}

\begin{frame}{Checklist for Choosing Controls}
  Before adding a variable \(X\) to a causal regression, ask:
  \begin{enumerate}
    \item Was \(X\) determined before treatment?
    \item Could treatment itself change \(X\)?
    \item Does conditioning on \(X\) make assignment more plausibly as-good-as-random?
    \item Do we still have overlap between treated and control units after conditioning?
  \end{enumerate}
  \begin{block}{Student habit to build}
    Do not say ``I controlled for everything available.'' Say exactly why each control belongs in the design.
  \end{block}
\end{frame}

\section{Matching and Propensity Scores}

\begin{frame}{Heterogeneity, Matching, and Propensity Scores}
  \begin{itemize}
    \item Constant treatment effects and linear CEFs are convenient but not necessary.
    \item Under CIA, matching and regression solve the same comparability problem.
    \item The main differences are how they average cell-specific effects and how much modeling they impose.
  \end{itemize}
\end{frame}

\begin{frame}{Why Heterogeneity Changes the Discussion}
  \begin{itemize}
    \item People may respond differently to the same treatment.
    \item The causal CEF may be nonlinear even if the average effect is well-defined.
    \item Regression then becomes a weighted average summary of many local comparisons.
  \end{itemize}
  \begin{block}{Key insight}
    Instead of abandoning regression, reinterpret it as a computationally attractive matching estimator.
  \end{block}
\end{frame}

\begin{frame}{Matching the Effect of Treatment on the Treated}
  \hypertarget{main:matching}{}
  With binary treatment \(D_i\), define
  \[
    \delta_X \equiv \E[Y_i \mid X_i,D_i=1]-\E[Y_i \mid X_i,D_i=0].
  \]
  Under CIA,
  \[
    \TOT
    =
    \E[Y_{1i}-Y_{0i}\mid D_i=1]
    =
    \E[\delta_X \mid D_i=1].
  \]
  If \(X_i\) is discrete,
  \[
    \TOT = \sum_x \delta_x \Prob(X_i=x \mid D_i=1).
  \]
  \begin{itemize}
    \item Matching constructs \(\delta_x\) cell by cell.
    \item Then it averages those cell effects using treated-group weights.
  \end{itemize}
  \derivbutton{matching-proof}
\end{frame}

\begin{frame}{ATT and ATE Differ by Weights}
  Under CIA,
  \[
    \ATE = \E[\delta_X] = \sum_x \delta_x \Prob(X_i=x),
  \]
  while
  \[
    \TOT = \sum_x \delta_x \Prob(X_i=x \mid D_i=1).
  \]
  \begin{itemize}
    \item Same covariate-specific treatment effects.
    \item Different weighting schemes.
    \item Therefore \(\ATE\) and \(\TOT\) need not be numerically close.
  \end{itemize}
  \begin{block}{Interpretation}
    \(\ATE\) answers ``what would treatment do for the population?'' while \(\TOT\) answers ``what did treatment do for those actually treated?''
  \end{block}
\end{frame}

\begin{frame}{Saturated Regression as a Matching Device}
  Consider the regression
  \[
    Y_i=\sum_x d_{ix}\beta_x+\delta_R D_i+\varepsilon_i,
  \]
  where \(d_{ix}\) is a dummy for the cell \(X_i=x\).
  \begin{itemize}
    \item This regression saturates the model in \(X_i\).
    \item It uses the same CIA logic as exact covariate matching.
    \item The coefficient \(\delta_R\) is not arbitrary; it is a weighted average of the same cell-specific effects \(\delta_x\).
  \end{itemize}
  \derivbutton{reg-matching-proof}
\end{frame}

\begin{frame}{Regression Is Variance-Weighted Matching}
  \hypertarget{main:reg-weight}{}
  A useful result is that
  \[
    \delta_R
    =
    \frac{\E[\sigma_D^2(X_i)\delta_X]}{\E[\sigma_D^2(X_i)]},
  \]
  where
  \[
    \sigma_D^2(X_i)=\Var(D_i \mid X_i)=p(X_i)(1-p(X_i)).
  \]
  \begin{itemize}
    \item Regression puts more weight on cells where treatment status varies a lot.
    \item Those are cells with both treated and control units, especially when treatment probability is near one half.
    \item Matching for \(\TOT\) instead emphasizes cells where treated units are common.
  \end{itemize}
  \derivbutton{reg-matching-proof}
\end{frame}

\begin{frame}{Why Different Weights Matter}
  \begin{itemize}
    \item If treatment effects are constant across covariate cells, all reasonable averages coincide.
    \item If treatment effects vary with \(X_i\), weighting matters.
    \item In the military-service application, those most likely to serve seem to benefit least, so \(\TOT\) matching estimates are smaller than regression estimates.
  \end{itemize}
  \begin{block}{Exam takeaway}
    ``Regression and matching disagree'' does not automatically mean one is wrong. They may be targeting different weighted averages.
  \end{block}
\end{frame}

\begin{frame}{Common Support}
  \begin{itemize}
    \item For causal comparisons at covariate value \(x\), we need both treated and untreated observations there.
    \item Formally, for binary treatment we want
    \[
      0<\Prob(D_i=1 \mid X_i=x)<1.
    \]
    \item If a cell contains only treated or only untreated units, \(\delta_x\) is undefined.
  \end{itemize}
  \begin{block}{Good news}
    In theory, both matching and saturated regression put zero weight on unsupported cells.
  \end{block}
  \begin{block}{Bad news}
    In practice, unsaturated parametric regressions can extrapolate into unsupported regions.
  \end{block}
\end{frame}

\begin{frame}{The Propensity Score Theorem}
  \hypertarget{main:pscore}{}
  Define the propensity score
  \[
    p(X_i)\equiv \E[D_i \mid X_i]=\Prob(D_i=1 \mid X_i).
  \]
  Rosenbaum and Rubin's theorem says:
  \begin{block}{Propensity-score theorem}
    If \(Y_{0i},Y_{1i}\indep D_i \mid X_i\), then \(Y_{0i},Y_{1i}\indep D_i \mid p(X_i)\).
  \end{block}
  \begin{itemize}
    \item A high-dimensional covariate vector can be replaced by a scalar balancing score.
    \item This is dimension reduction, not magic.
  \end{itemize}
  \derivbutton{pscore-proof}
\end{frame}

\begin{frame}{How the Propensity Score Is Used}
  Under CIA,
  \[
    \TOT
    =
    \E\left\{
      \E[Y_i \mid p(X_i),D_i=1]-\E[Y_i \mid p(X_i),D_i=0]
      \mid D_i=1
    \right\}.
  \]
  Practical implementations:
  \begin{itemize}
    \item match treated and control units with similar estimated scores,
    \item stratify observations into score bins,
    \item weight observations by inverse probabilities,
    \item screen observations to impose common support before regression.
  \end{itemize}
\end{frame}

\begin{frame}{Inverse-Probability Weighting Identities}
  \hypertarget{main:ipw}{}
  Under CIA,
  \[
    \ATE
    =
    \E\left[
      \frac{Y_iD_i}{p(X_i)}
      -
      \frac{Y_i(1-D_i)}{1-p(X_i)}
    \right].
  \]
  Equivalently,
  \[
    \ATE
    =
    \E\left[
      \frac{(D_i-p(X_i))Y_i}{p(X_i)(1-p(X_i))}
    \right].
  \]
  For the treated effect,
  \[
    \TOT
    =
    \E\left[
      \frac{(D_i-p(X_i))Y_i}{(1-p(X_i))\Prob(D_i=1)}
    \right].
  \]
  \derivbutton{ipw-proof}
\end{frame}

\begin{frame}{Regression, Matching, and IPW Are One Family}
  A unified way to write these estimands is:
  \[
    \E\left\{
      g(X_i)\left[
        \frac{Y_iD_i}{p(X_i)}
        -
        \frac{Y_i(1-D_i)}{1-p(X_i)}
      \right]
    \right\}.
  \]
  \begin{itemize}
    \item \(g(X_i)=1\) gives the ATE.
    \item \(g(X_i)=p(X_i)/\Prob(D_i=1)\) gives the treated effect.
    \item A regression-type choice of \(g(X_i)\) reproduces the regression estimand.
  \end{itemize}
  \begin{block}{Bottom line}
    These methods differ mostly in how they weight the same underlying cell-level comparisons.
  \end{block}
\end{frame}

\begin{frame}{What to Conclude About Propensity Scores}
  \begin{itemize}
    \item Propensity scores are attractive when treatment assignment is easier to model than outcomes.
    \item They help reduce dimensionality and organize overlap diagnostics.
    \item But they are not automatically more causal or more robust than regression.
    \item Regression with flexible controls often performs very well.
  \end{itemize}
  \begin{block}{Practical recommendation}
    Start with regression. Use propensity-score tools as complements, especially for overlap screening and dimension reduction.
  \end{block}
\end{frame}

\begin{frame}{NSW Lesson: Design Beats Method Branding}
  \begin{itemize}
    \item Lalonde showed naive observational comparisons can fail badly.
    \item Dehejia and Wahba showed better overlap and better covariate control can move observational estimates toward the experimental benchmark.
    \item Careful regression control plus support screening can be as persuasive as propensity-score matching.
  \end{itemize}
  \begin{block}{Applied lesson}
    Do not ask ``Should I use regression or propensity scores?'' Ask ``What makes CIA plausible here, and where is the common support?''
  \end{block}
\end{frame}

\section{Practical Issues}

\begin{frame}{Practical Issues That Cause Confusion}
  \begin{itemize}
    \item Sample weights
    \item Grouped-data regressions
    \item Limited dependent variables
    \item Conditional-on-positive comparisons
  \end{itemize}
  \begin{block}{Theme}
    Many practical habits are defensible only if they move us closer to the population estimand we actually care about.
  \end{block}
\end{frame}

\begin{frame}{When Should You Weight a Regression?}
  \begin{itemize}
    \item Use sampling weights when the sample is nonrandom and weights are inverse sampling probabilities.
    \item Then weighted least squares targets the population regression function rather than the sample-specific one.
    \item Use group-size weights when regressing cell means and you want to recover the microdata regression.
  \end{itemize}
  \begin{block}{Rule of thumb}
    Weight when weighting makes your sample look more like the population target.
  \end{block}
\end{frame}

\begin{frame}{Grouped Data and Weighting by Cell Size}
  \hypertarget{main:group-weight}{}
  Suppose you observe cell means \(\bar{y}_x=\E[Y_i \mid X_i=x]\) and cell frequencies \(n_x\).
  \begin{itemize}
    \item Regressing \(\bar{y}_x\) on \(x\) with weights \(n_x\) reproduces the microdata regression.
    \item Unweighted regression on cell means changes the estimand because small and large cells get equal influence.
  \end{itemize}
  \begin{block}{Interpretation}
    Weighting by cell size is not a cosmetic choice; it determines whether we are approximating the population or a fictional equally-weighted set of cells.
  \end{block}
  \derivbutton{grouped-proof}
\end{frame}

\begin{frame}{Why We Are Skeptical of Weighting for Heteroskedasticity}
  \begin{itemize}
    \item Heteroskedasticity does not change the population regression coefficient.
    \item Poorly estimated variance weights can hurt finite-sample performance.
    \item If the CEF is nonlinear, weighted least squares may approximate a different object from OLS.
  \end{itemize}
  \begin{block}{Practical advice}
    Usually keep OLS for the point estimate and use heteroskedasticity-robust standard errors for inference.
  \end{block}
\end{frame}

\begin{frame}{Limited Dependent Variables Do Not Kill OLS}
  \begin{itemize}
    \item If treatment is randomly assigned, the coefficient on a treatment dummy still equals the difference in means.
    \item This remains true whether \(Y_i\) is continuous, binary, or nonnegative.
    \item For a binary outcome, the difference in means is also a difference in probabilities.
  \end{itemize}
  \[
    \E[Y_i \mid D_i=1]-\E[Y_i \mid D_i=0]=\E[Y_{1i}-Y_{0i}].
  \]
  \begin{block}{Message}
    The outcome being discrete or bounded changes interpretation, not the validity of the mean comparison itself.
  \end{block}
\end{frame}

\begin{frame}{Why People Use Probit and Tobit Anyway}
  \begin{itemize}
    \item Nonlinear models can fit bounded outcomes more naturally.
    \item Probit fitted values stay between zero and one.
    \item Tobit fitted values stay nonnegative under the latent model.
  \end{itemize}
  \begin{block}{But}
    The raw coefficients in Probit or Tobit are not the causal effects on observed outcomes. They must be converted into finite differences or average derivatives.
  \end{block}
\end{frame}

\begin{frame}{Marginal Effects Are the Right Comparison Object}
  For a nonlinear model with covariates, compare OLS to
  \[
    \E\left\{\E[Y_i \mid X_i,D_i=1]-\E[Y_i \mid X_i,D_i=0]\right\},
  \]
  not to latent-index coefficients.
  \begin{itemize}
    \item OLS estimates an average effect on the observed outcome.
    \item Probit and Tobit coefficients live on latent scales.
    \item Once converted to marginal effects, nonlinear and OLS estimates are often quite similar in practice.
  \end{itemize}
\end{frame}

\begin{frame}{Conditional-on-Positive Effects Are Another Bad-Control Problem}
  \hypertarget{main:cop}{}
  For a nonnegative outcome,
  \[
    \E[Y_i \mid D_i]
    =
    \E[Y_i \mid Y_i>0,D_i]\Prob(Y_i>0 \mid D_i).
  \]
  Some researchers then compare
  \[
    \E[Y_i \mid Y_i>0,D_i=1]-\E[Y_i \mid Y_i>0,D_i=0].
  \]
  \begin{itemize}
    \item This conditions on an outcome that treatment may itself change.
    \item So it is a bad-control comparison, even in a randomized experiment.
  \end{itemize}
  \derivbutton{cop-proof}
\end{frame}

\begin{frame}{Why COP Effects Are Not Causal}
  \begin{itemize}
    \item Treatment changes who is in the \(Y_i>0\) group.
    \item Therefore treated positives and control positives are different populations.
    \item Even if the true treatment effect were zero, the conditional-on-positive comparison could still be nonzero.
  \end{itemize}
  \begin{block}{Safe object}
    The average effect on the observed outcome.
  \end{block}
  \begin{block}{Unsafe object}
    The effect among those with positive observed outcomes after treatment assignment has already changed that set.
  \end{block}
\end{frame}

\begin{frame}{When Tobit-Type Models Do Make Sense}
  \begin{itemize}
    \item Tobit is most defensible under true censoring.
    \item Example: earnings are top-coded in a survey, but the latent uncensored earnings are the object of interest.
    \item It is much less convincing when zeros are genuine outcomes, such as zero medical spending or zero hours worked.
  \end{itemize}
  \begin{block}{Practical position}
    If the latent variable has no meaningful empirical interpretation, do not hide behind it.
  \end{block}
\end{frame}

\begin{frame}{Practical Modeling Rules for Students}
  \begin{enumerate}
    \item Start with the average effect on the observed outcome.
    \item Use flexible pre-treatment controls and robust standard errors.
    \item Check overlap before trusting regression or matching.
    \item Treat Probit/Tobit as optional curve-fitting devices, not automatic upgrades.
    \item Never condition on variables created by treatment unless you have a very explicit structural reason.
  \end{enumerate}
\end{frame}

\section{Synthesis}

\begin{frame}{Integrated Summary of Chapter 3}
  \begin{itemize}
    \item The CEF is the fundamental population object; regression is its best linear approximation.
    \item A regression becomes causal only when the underlying CEF is causal, typically under CIA plus good controls.
    \item OVB explains mechanically why coefficients move when controls are added or omitted.
    \item Matching, propensity-score weighting, and regression all average cell-specific causal effects using different weights.
    \item Limited dependent variables change presentation and modeling choices, but not the centrality of average effects.
  \end{itemize}
\end{frame}

\begin{frame}{What Students Should Be Able to Do After This Lecture}
  \begin{itemize}
    \item Define the CEF and explain why regression approximates it.
    \item State the conditional independence assumption in words and symbols.
    \item Use the omitted variables bias formula to reason about coefficient movements.
    \item Explain why occupation can be a bad control in a schooling regression.
    \item Distinguish \(\ATE\), \(\TOT\), and regression-weighted effects.
    \item Explain what the propensity score does and does not do.
    \item Defend when to use weights, and when not to.
    \item Explain why conditional-on-positive comparisons are not generally causal.
  \end{itemize}
\end{frame}

\appendix

\begin{frame}{Appendix Map}
  \begin{itemize}
    \item A1: Law of iterated expectations and CEF decomposition
    \item A2: Why regression approximates the CEF
    \item A3: Regression anatomy
    \item A4: CIA and the linear causal regression
    \item A5: Omitted variables bias
    \item A6: Bad control decomposition
    \item A7: Matching formulas for \(\TOT\) and \(\ATE\)
    \item A8: Regression as variance-weighted matching
    \item A9: Propensity-score theorem
    \item A10: Inverse-probability weighting identities
    \item A11: Grouped-data weighting
    \item A12: Conditional-on-positive decomposition
  \end{itemize}
\end{frame}

\begin{frame}{Appendix A1: LIE and CEF Decomposition}
  \hypertarget{app:cef-foundations}{}
  Start from the law of iterated expectations:
  \[
    \E[Y_i]=\E\{\E[Y_i \mid X_i]\}.
  \]
  Define
  \[
    \varepsilon_i \equiv Y_i-\E[Y_i \mid X_i].
  \]
  Then
  \[
    \E[\varepsilon_i \mid X_i]
    =
    \E[Y_i \mid X_i]-\E[Y_i \mid X_i]
    =
    0.
  \]
  Therefore
  \[
    Y_i=\E[Y_i \mid X_i]+\varepsilon_i
    \quad \text{with} \quad
    \E[\varepsilon_i \mid X_i]=0.
  \]
  Any function \(h(X_i)\) is also uncorrelated with \(\varepsilon_i\) because
  \[
    \E[h(X_i)\varepsilon_i]
    =
    \E\{h(X_i)\E[\varepsilon_i \mid X_i]\}
    =
    0.
  \]
  \backbutton{lie}
\end{frame}

\begin{frame}{Appendix A2: Why Regression Approximates the CEF}
  \hypertarget{app:reg-justification}{}
  Write
  \[
    \left(Y_i-X_i'b\right)^2
    =
    \left(Y_i-\E[Y_i \mid X_i]\right)^2
    +
    \left(\E[Y_i \mid X_i]-X_i'b\right)^2
    +
    2\left(Y_i-\E[Y_i \mid X_i]\right)\left(\E[Y_i \mid X_i]-X_i'b\right).
  \]
  Take expectations.
  \begin{itemize}
    \item The first term does not depend on \(b\).
    \item The cross term has expectation zero because \(Y_i-\E[Y_i \mid X_i]\) is orthogonal to any function of \(X_i\).
  \end{itemize}
  So minimizing \(\E[(Y_i-X_i'b)^2]\) is equivalent to minimizing
  \[
    \E\left[\left(\E[Y_i \mid X_i]-X_i'b\right)^2\right].
  \]
  Hence regression is the best linear approximation to the CEF.
  \backbutton{reg-just}
\end{frame}

\begin{frame}{Appendix A3: Regression Anatomy}
  \hypertarget{app:anatomy-proof}{}
  Let \(\tilde{x}_{ki}\) be the residual from regressing \(x_{ki}\) on all other covariates.
  Start from
  \[
    Y_i=\beta_0+\beta_k x_{ki}+\text{other terms}+e_i.
  \]
  Take covariance with \(\tilde{x}_{ki}\):
  \[
    \Cov(Y_i,\tilde{x}_{ki})
    =
    \beta_k \Cov(x_{ki},\tilde{x}_{ki}),
  \]
  because
  \begin{itemize}
    \item \(\tilde{x}_{ki}\) is uncorrelated with all other regressors by construction,
    \item \(\tilde{x}_{ki}\) is uncorrelated with the regression residual.
  \end{itemize}
  Since \(\Cov(x_{ki},\tilde{x}_{ki})=\Var(\tilde{x}_{ki})\),
  \[
    \beta_k=\frac{\Cov(Y_i,\tilde{x}_{ki})}{\Var(\tilde{x}_{ki})}.
  \]
  \backbutton{anatomy}
\end{frame}

\begin{frame}{Appendix A4: CIA and the Linear Causal Regression}
  \hypertarget{app:cia-proof}{}
  With binary treatment, CIA says
  \[
    \{Y_{0i},Y_{1i}\}\indep C_i \mid X_i.
  \]
  Therefore
  \[
    \E[Y_{0i} \mid X_i,C_i=1]=\E[Y_{0i} \mid X_i,C_i=0]
  \]
  and similarly for \(Y_{1i}\). Hence
  \[
    \E[Y_i \mid X_i,C_i=1]-\E[Y_i \mid X_i,C_i=0]
    =
    \E[Y_{1i}-Y_{0i} \mid X_i].
  \]
  For the linear causal model
  \[
    f_i(s)=\alpha+\rho s+\eta_i,
    \qquad
    \eta_i=X_i'\gamma+v_i,
  \]
  CIA implies
  \[
    \E[f_i(s)\mid X_i,S_i]=\E[f_i(s)\mid X_i]=\alpha+\rho s+X_i'\gamma,
  \]
  so the observed-outcome equation
  \[
    Y_i=\alpha+\rho S_i+X_i'\gamma+v_i
  \]
  has a residual orthogonal to \(S_i\) and \(X_i\), and \(\rho\) is causal.
  \backbutton{cia}
\end{frame}

\begin{frame}{Appendix A5: Omitted Variables Bias}
  \hypertarget{app:ovb-proof}{}
  Long regression:
  \[
    Y_i=\alpha+\rho S_i+A_i'\gamma+\varepsilon_i.
  \]
  The short-regression slope on \(S_i\) is
  \[
    \frac{\Cov(Y_i,S_i)}{\Var(S_i)}.
  \]
  Substitute the long regression into the covariance:
  \[
    \Cov(Y_i,S_i)
    =
    \rho \Var(S_i)+\Cov(A_i'\gamma,S_i)+\Cov(\varepsilon_i,S_i).
  \]
  The last term is zero in the long regression, so
  \[
    \frac{\Cov(Y_i,S_i)}{\Var(S_i)}
    =
    \rho+\frac{\Cov(A_i'\gamma,S_i)}{\Var(S_i)}.
  \]
  Writing \(\delta_{AS}\) for the coefficients from regressions of \(A_i\) on \(S_i\),
  \[
    \frac{\Cov(Y_i,S_i)}{\Var(S_i)}
    =
    \rho+\gamma'\delta_{AS}.
  \]
  \backbutton{ovb}
\end{frame}

\begin{frame}{Appendix A6: Why Bad Controls Generate Selection Bias}
  \hypertarget{app:bad-control-proof}{}
  Let \(C_i\) be college and \(W_i\) be white-collar status, with potential outcomes
  \[
    Y_i=C_iY_{1i}+(1-C_i)Y_{0i},
    \qquad
    W_i=C_iW_{1i}+(1-C_i)W_{0i}.
  \]
  Even if \(C_i\) is randomly assigned,
  \[
    \E[Y_i \mid W_i=1,C_i=1]-\E[Y_i \mid W_i=1,C_i=0]
  \]
  becomes
  \[
    \E[Y_{1i} \mid W_{1i}=1]-\E[Y_{0i} \mid W_{0i}=1].
  \]
  Add and subtract \(\E[Y_{0i} \mid W_{1i}=1]\):
  \[
    \E[Y_{1i}-Y_{0i} \mid W_{1i}=1]
    +
    \left\{
      \E[Y_{0i} \mid W_{1i}=1]-\E[Y_{0i} \mid W_{0i}=1]
    \right\}.
  \]
  The second term is selection bias caused by conditioning on a post-treatment variable.
  \backbutton{bad-control}
\end{frame}

\begin{frame}{Appendix A7: Matching Formulas}
  \hypertarget{app:matching-proof}{}
  Define
  \[
    \delta_X\equiv \E[Y_i \mid X_i,D_i=1]-\E[Y_i \mid X_i,D_i=0].
  \]
  Under CIA,
  \[
    \E[Y_{0i}\mid X_i,D_i=1]=\E[Y_{0i}\mid X_i,D_i=0],
  \]
  so
  \[
    \TOT
    =
    \E[Y_{1i}-Y_{0i}\mid D_i=1]
    =
    \E[\delta_X \mid D_i=1].
  \]
  If \(X_i\) is discrete,
  \[
    \TOT=\sum_x \delta_x \Prob(X_i=x \mid D_i=1).
  \]
  Similarly,
  \[
    \ATE=\E[\delta_X]=\sum_x \delta_x \Prob(X_i=x).
  \]
  \backbutton{matching}
\end{frame}

\begin{frame}{Appendix A8: Regression as Variance-Weighted Matching}
  \hypertarget{app:reg-matching-proof}{}
  With saturated controls in \(X_i\), the regression coefficient on \(D_i\) satisfies
  \[
    \delta_R
    =
    \frac{\Cov(Y_i,\tilde{D}_i)}{\Var(\tilde{D}_i)}
    =
    \frac{\E[(D_i-\E[D_i \mid X_i])Y_i]}{\E[(D_i-\E[D_i \mid X_i])^2]}.
  \]
  Replace \(Y_i\) by \(\E[Y_i \mid D_i,X_i]\) and expand the CEF:
  \[
    \E[Y_i \mid D_i,X_i]
    =
    \E[Y_i \mid D_i=0,X_i]+\delta_X D_i.
  \]
  The baseline term drops out because it is a function of \(X_i\). What remains is
  \[
    \delta_R
    =
    \frac{\E[(D_i-\E[D_i \mid X_i])^2\delta_X]}{\E[(D_i-\E[D_i \mid X_i])^2]}
    =
    \frac{\E[\sigma_D^2(X_i)\delta_X]}{\E[\sigma_D^2(X_i)]}.
  \]
  For binary treatment, \(\sigma_D^2(X_i)=p(X_i)(1-p(X_i))\).
  \backbutton{reg-weight}
\end{frame}

\begin{frame}{Appendix A9: Proof of the Propensity-Score Theorem}
  \hypertarget{app:pscore-proof}{}
  We want to show that \(Y_{ji}\indep D_i \mid p(X_i)\), where \(j\in\{0,1\}\).
  It is enough to prove that
  \[
    \Prob(D_i=1 \mid Y_{ji},p(X_i))
  \]
  does not depend on \(Y_{ji}\).
  Using iterated expectations,
  \[
  \begin{aligned}
    \Prob(D_i=1 \mid Y_{ji},p(X_i))
    &= \E[D_i \mid Y_{ji},p(X_i)] \\
    &= \E\{\E[D_i \mid Y_{ji},p(X_i),X_i] \mid Y_{ji},p(X_i)\} \\
    &= \E\{\E[D_i \mid Y_{ji},X_i] \mid Y_{ji},p(X_i)\} \\
    &= \E\{\E[D_i \mid X_i] \mid Y_{ji},p(X_i)\},
  \end{aligned}
  \]
  where the last step uses CIA given \(X_i\).
  But \(\E[D_i \mid X_i]=p(X_i)\), so the final expression is simply \(p(X_i)\), which does not depend on \(Y_{ji}\).
  \backbutton{pscore}
\end{frame}

\begin{frame}{Appendix A10: Inverse-Probability Weighting Identities}
  \hypertarget{app:ipw-proof}{}
  Under CIA,
  \[
    \E\left[\frac{Y_iD_i}{p(X_i)}\right]
    =
    \E\left\{
      \E\left[\frac{Y_iD_i}{p(X_i)} \middle| X_i \right]
    \right\}
    =
    \E[Y_{1i}].
  \]
  Similarly,
  \[
    \E\left[\frac{Y_i(1-D_i)}{1-p(X_i)}\right]
    =
    \E[Y_{0i}].
  \]
  Subtract to obtain
  \[
    \ATE
    =
    \E\left[
      \frac{Y_iD_i}{p(X_i)}-\frac{Y_i(1-D_i)}{1-p(X_i)}
    \right].
  \]
  Rearranging terms yields
  \[
    \ATE
    =
    \E\left[
      \frac{(D_i-p(X_i))Y_i}{p(X_i)(1-p(X_i))}
    \right].
  \]
  \backbutton{ipw}
\end{frame}

\begin{frame}{Appendix A11: Why Group-Size Weighting Recovers the Microdata Regression}
  \hypertarget{app:grouped-proof}{}
  Suppose \(X_i\) takes discrete values \(x\), and for each cell you observe
  \[
    \bar{y}_x = \E[Y_i \mid X_i=x]
    \qquad \text{and} \qquad
    n_x.
  \]
  The population regression moment is
  \[
    \E[X_iY_i]
    =
    \sum_x x\,\E[Y_i \mid X_i=x]\Prob(X_i=x).
  \]
  In a sample, \(\Prob(X_i=x)\) is estimated by \(n_x/N\), so
  \[
    \widehat{\E[X_iY_i]}
    =
    \sum_x x\,\bar{y}_x \frac{n_x}{N}.
  \]
  This is exactly the cross-moment generated by a weighted regression of \(\bar{y}_x\) on \(x\) with weights \(n_x\).
  Therefore weighting cell means by cell size reproduces the microdata regression target.
  \backbutton{group-weight}
\end{frame}

\begin{frame}{Appendix A12: Conditional-on-Positive Decomposition}
  \hypertarget{app:cop-proof}{}
  Consider the conditional-on-positive comparison:
  \[
    \E[Y_i \mid Y_i>0,D_i=1]-\E[Y_i \mid Y_i>0,D_i=0].
  \]
  In a randomized experiment this equals
  \[
    \E[Y_{1i} \mid Y_{1i}>0]-\E[Y_{0i} \mid Y_{0i}>0].
  \]
  Add and subtract \(\E[Y_{0i} \mid Y_{1i}>0]\):
  \[
  \begin{aligned}
    &\E[Y_{1i}-Y_{0i} \mid Y_{1i}>0] \\
    &\quad + \left\{
      \E[Y_{0i} \mid Y_{1i}>0]-\E[Y_{0i} \mid Y_{0i}>0]
    \right\}.
  \end{aligned}
  \]
  The second term is selection bias:
  treatment changes who ends up in the positive-outcome group, so the two groups being compared are not the same population.
  \backbutton{cop}
\end{frame}

\end{document}
